\documentclass[11pt,a4paper]{article}
\usepackage[utf8]{inputenc}
\usepackage[T1]{fontenc}
\usepackage[brazilian]{babel}
\usepackage[margin=2.3cm]{geometry}
\usepackage{booktabs,array,tabularx}
\usepackage{xcolor}
\usepackage{pgfplots}
\pgfplotsset{compat=1.17}
\PassOptionsToPackage{hyphens}{url}
\usepackage[hidelinks]{hyperref}
\usepackage{enumitem}
\setlist{nosep,leftmargin=1.4em}
\emergencystretch=2em

\definecolor{sol}{HTML}{3B6EA8}
\definecolor{perda}{HTML}{C0504D}
\definecolor{pergunta}{HTML}{E0A030}

\title{Sweatshop no PhysiSyst: avisos do simulador, recusa no próprio corpo e corda solta da polia ideal\\\large GPT-6.1 Sol (high) no stage 2, Sonnet 5.5 (xhigh) no hard, GPT-6.1 Sol (max) no stage 3}
\author{Foreman, sessão \texttt{sweatshop/2026-10-01-2342}}
\date{2 de outubro de 2026}

\begin{document}
\maketitle

\section*{Resumo}

Run noturno com
\texttt{-Models 'stage 2 gpt-6.1-sol high, hard sonnet-5.5 xhigh, stage 3
gpt-6.1-sol max'} (lineup \emph{ad hoc}, o \texttt{Models (sol)} com o hard
subido de high para xhigh), a pedido do Bruno: ``seguir com os tickets
restantes''. A fila eram CLEAN-24, PHY-51, PHY-52 e PHY-54 (\texttt{Difficulty:
hard}), com o PHY-53 atrás do PHY-51.

\begin{itemize}
  \item \textbf{Produzido:} 5 tickets mergeados, com \textbf{uma reabertura}
        (PHY-54, por fronteira de Primary files). PR de sessão:
        \href{https://github.com/Laginho/PhysiSyst/pull/12}{\#12}, 5 linhas,
        todas \texttt{Review: agent} --- Approve.
  \item \textbf{Custo:} \textbf{\$15,32} a preço de lista em 12 dos 13
        estágios registrados: \texttt{gpt-6.1-sol high} \$2,50 (4),
        \texttt{sonnet-5.5 xhigh} \$8,82 (2), \texttt{gpt-6.1-sol max} \$4,01
        (6); uma linha saiu com \texttt{?}. 128 min de estágio, 29,8\,M tokens
        de entrada (95,2\,\% em cache) e 250\,k de saída. \$3,06 por ticket
        mergeado.
  \item \textbf{O que deu errado:} o limite de uso do Codex. Às 00:45, 3 min
        dentro da revisão do PHY-53, o Codex respondeu ``try again at
        3:49 AM''. O driver esperou até 03:51, como deveria, mas o run ficou
        3\,h\,06 parado: 58\,\% das 5\,h\,18 de relógio.
  \item \textbf{Veredito provisório:} o Sol high entregou 4 de 4 em
        \texttt{to-review}, sem reabrir, a \$0,62 por implementação. O Sonnet
        5.5 xhigh resolveu o ticket mais difícil da fila e acertou a física de
        primeira, mas custou \$8,82 por um ticket e reabriu por escopo. Com 2
        revisões, é anedota.
\end{itemize}

\section{Linha do tempo}

\begin{tabularx}{\linewidth}{@{}lX@{}}
\toprule
23:42 & Lançamento único, driver \texttt{74dc512}. Sessão nova
        \texttt{sweatshop/2026-10-01-2342} (o PR \#11 já estava mergeado).\\
23:52 & CLEAN-24 (recusa visível no clique no próprio corpo) em
        \texttt{to-review} em 9 min. 00:01: mergeado, sem achados.\\
00:09 & PHY-51 (aviso do teto de $\omega$ do Rapier na API) em 8 min. 00:16:
        mergeado; a revisão só ajustou a ordem de comentários.\\
00:23 & PHY-52 (residual zerado além do \texttt{ROPE\_SLIP} na polia com
        massa) em 6 min. 00:32: mergeado, sem achados.\\
00:41 & PHY-53 (avisos do simulador na tela) em 10 min.\\
00:45 & A revisão do PHY-53 bate no limite de uso do Codex aos 3 min; o
        driver registra \texttt{api error, not counted} e espera até 03:51.\\
03:51 & Revisão do PHY-53 recomeça. 04:03: mergeado, sem achados.\\
04:03 & PHY-54 vai para o implementador hard (\texttt{sonnet-5.5 xhigh}). Em
        19 min, \texttt{to-review}: a corda se solta da polia ideal em vez de
        enrolar uma volta.\\
04:40 & \textbf{PHY-54 reaberto} em 18 min: os critérios passam, mas a guarda
        de \texttt{pieceLengths} e o bloco de testes estão fora dos Primary
        files. A revisão abre o CLEAN-25.\\
04:47 & Passada 2 em 6 min. O próprio stage 2 chama o proxy, que amplia o
        Primary e cria o critério 6; os testes mudam de bloco num commit só de
        teste.\\
04:59 & PHY-54 mergeado em 13 min. 05:00: ``Nothing runnable''; PR \#12
        aberto.\\
\bottomrule
\end{tabularx}

\section{Métricas por estágio}

Linhas 116--128 de \texttt{.scratch/run-log.md}. Minutos, tokens e custo como o
driver gravou.

\medskip
\noindent{\footnotesize
\begin{tabularx}{\linewidth}{@{}l l r r r r r X@{}}
\toprule
Estágio & Modelo & Min & Entrada & Cache & Saída & Custo & Desfecho\\
\midrule
C24 impl.      & sol high     & 9  & 2,40\,M & 97\,\% & 8\,k   & 0,694 & to-review\\
C24 revisão    & sol max      & 9  & 1,70\,M & 93\,\% & 12\,k  & 0,649 & mergeado\\
P51 impl.      & sol high     & 8  & 1,70\,M & 96\,\% & 8\,k   & 0,535 & to-review\\
P51 revisão    & sol max      & 6  & 0,89\,M & 93\,\% & 9\,k   & 0,385 & mergeado\\
P52 impl.      & sol high     & 6  & 1,20\,M & 90\,\% & 8\,k   & 0,560 & to-review\\
P52 revisão    & sol max      & 9  & 1,80\,M & 95\,\% & 13\,k  & 0,652 & mergeado\\
P53 impl.      & sol high     & 10 & 2,20\,M & 96\,\% & 11\,k  & 0,707 & to-review\\
P53 revisão    & sol max      & 3  & ?       &        & ?      & ?     & limite de uso\\
P53 revisão    & sol max      & 12 & 1,80\,M & 96\,\% & 13\,k  & 0,628 & mergeado\\
P54 impl.      & sonnet xhigh & 19 & 8,10\,M & 96\,\% & 110\,k & 6,058 & to-review\\
P54 revisão    & sol max      & 18 & 2,90\,M & 96\,\% & 20\,k  & 0,998 & reaberto\\
P54 impl.\ 2   & sonnet xhigh & 6  & 3,30\,M & 95\,\% & 25\,k  & 2,757 & to-review\\
P54 revisão 2  & sol max      & 13 & 1,80\,M & 93\,\% & 13\,k  & 0,693 & mergeado\\
\midrule
\multicolumn{2}{@{}l}{Total} & 128 & 29,79\,M & 95,2\,\% & 250\,k & 15,316 & 12 de 13 com preço\\
\multicolumn{2}{@{}l}{Implementação} & 58 & 18,90\,M & 95,6\,\% & 170\,k & 11,311 & 6 estágios\\
\multicolumn{2}{@{}l}{Revisão} & 70 & 10,89\,M & 94,6\,\% & 80\,k & 4,005 & 6 estágios com preço\\
\multicolumn{2}{@{}l}{Retrabalho} & 19 & 5,10\,M & & 38\,k & 3,450 & 2 estágios (reabertura)\\
\multicolumn{2}{@{}l}{Perdido} & 3 & ? & & ? & ? & 1 linha sem preço\\
\bottomrule
\end{tabularx}}

\bigskip
\begin{center}
\begin{tikzpicture}
\begin{axis}[
  ybar stacked, bar width=18pt, width=\linewidth, height=6.5cm,
  ylabel={Tokens de entrada (M)}, ymin=0,
  symbolic x coords={CLEAN-24,PHY-51,PHY-52,PHY-53,PHY-54},
  xtick=data, x tick label style={font=\small},
  legend style={at={(0.02,0.97)},anchor=north west,draw=none,font=\small},
  axis lines*=left, ymajorgrids, grid style={gray!25}]
\addplot[fill=sol,draw=none] coordinates
  {(CLEAN-24,4.10) (PHY-51,2.59) (PHY-52,3.00) (PHY-53,4.00) (PHY-54,11.00)};
\addplot[fill=perda,draw=none] coordinates
  {(CLEAN-24,0) (PHY-51,0) (PHY-52,0) (PHY-53,0) (PHY-54,5.10)};
\legend{primeira passada e revisão, retrabalho da reabertura}
\end{axis}
\end{tikzpicture}
\end{center}

\paragraph{Onde foi o custo.} Por ticket: PHY-54 \$10,51 (56 min), CLEAN-24
\$1,34, PHY-53 \$1,34 (mais 3 min perdidos), PHY-52 \$1,21, PHY-51 \$0,92. O
PHY-54 levou 69\,\% do custo do run; só a primeira implementação (\$6,06,
110\,k de saída) custou mais que os outros quatro tickets juntos (\$4,81). O
Sol high cobra \$0,62 por implementação; o Sonnet 5.5 xhigh, \$4,41; o Sol max,
\$0,67 por revisão.

\paragraph{O que se perdeu.} Em tokens, quase nada: uma linha \texttt{?} de
3 min. Em relógio, 3\,h\,06 de espera pelo reset do Codex (00:45--03:51). O
limite chegou depois de 7 estágios Codex completos e 11,9\,M tokens de entrada
desde 23:42.

\section{Comparação}

Os runs anteriores do PhysiSyst com preço, todos com
\texttt{gpt-6.1-sol max} na revisão.

\medskip
\noindent{\small
\begin{tabularx}{\linewidth}{@{}X r r r r@{}}
\toprule
& 30/09 (1--5) & 30/09 (6) & 01/10 & Este run\\
\midrule
Implementador & sol high & sol high & sonnet-5.5 high & sol high + sonnet xhigh\\
Estágios registrados & 28 & 10 & 21 & 13\\
Tickets mergeados & 10 & 2 & 7 & 5\\
Reaberturas & 1 & 3 & 1 & 1\\
Implementação: falhou ou perguntou & 4 de 15 & 0 de 5 & 0 de 8 & 0 de 6\\
Mediana da implementação & 7 min & 5 min & 6 min & 8,5 min\\
Custo médio por implementação & \$0,52 & \$0,45 & \$1,74 & \$1,89\\
\quad só Sol high & \$0,52 & \$0,45 & --- & \$0,62\\
Mediana da revisão (com preço) & 10 min & 11 min & 12,5 min & 10,5 min\\
Custo médio por revisão & \$0,59 & \$0,62 & \$0,65 & \$0,67\\
Estágios perdidos (\texttt{?}) & 2 & 0 & 5 & 1\\
Custo registrado & \$14,20 & \$5,39 & \$19,13 & \$15,32\\
Custo por ticket mergeado & \$1,42 & \$2,69 & \$2,73 & \$3,06\\
\bottomrule
\end{tabularx}}

\medskip
Sem o PHY-54, o run custaria \$4,81 por 4 tickets (\$1,20 cada), o melhor
número do PhysiSyst até aqui. O Sol high ficou um pouco mais caro que em 30/09
(\$0,62 contra \$0,45--0,52) com tickets de UI e física um pouco maiores. O
gate passou em todos os estágios sem \texttt{VITEST\_MAX\_WORKERS=1}: o
CLEAN-23 (PR \#11) resolveu o teste lento que o relatório anterior apontava.

\section{Atribuição das reaberturas}

Uma reabertura no run: PHY-54, rodada 1, às 04:40. O revisor foi
\texttt{gpt-6.1-sol max}, diferente do foreman (Opus 5.5), então o foreman
classificou. Evidência: o ticket como o stage 2 o recebeu (\texttt{14e3373},
o mesmo da base da sessão), o diff revisado (\texttt{f2f6d62...497fc6f}) e a
revisão (\texttt{a552a99}). Linhas em \path{.scratch/reopen-attribution.md}.

\medskip
\noindent{\small
\begin{tabularx}{\linewidth}{@{}l c l X@{}}
\toprule
Ticket & Rodada & Rótulo & Achado e porquê\\
\midrule
PHY-54 & 1 & S1 & A guarda \texttt{Math.max(0, \dots)} em
  \texttt{pieceLengths} está fora dos Primary files de \texttt{simulator.ts},
  embora a própria revisão a julgue necessária. O stage 1 já sabia que
  \texttt{pieceLengths} soma a varredura negativa, mas mandou só a corda mista
  para o PHY-55 e esqueceu a corda ideal em grupo. O conserto pediu um
  \texttt{Proxy decided} que ampliou o Primary e criou o critério 6.\\
PHY-54 & 1 & S2-expl./código ? & Os cinco testes novos de \texttt{ropePath}
  estão num \texttt{describe} próprio, fora do bloco nomeado nos Primary
  files. O bloco está escrito no ticket, e mover não pediu decisão. Fica em
  dúvida porque é de lugar, não de comportamento, e beira o S3.\\
PHY-54 & 1 & S2-implícito & \texttt{wrapAngle} de \texttt{ropePath.ts}
  repete a fórmula de \texttt{simulator.ts}. Duplicação do próprio stage 2,
  que a sinalizou e deixou; não bloqueou.\\
PHY-54 & 1 & S1 & Duas polias soltas com projeções em ordem inversa à de
  \texttt{via} dão 21\,m e puxões $(\mp2, 0)$ em vez da perna reta de 11\,m.
  Caso fora da regra das várias polias que o proxy fixou; vira o CLEAN-25 e
  não bloqueou.\\
\bottomrule
\end{tabularx}}

\medskip
Totais: 2 S1, 1 S2-explícito/código (em dúvida), 1 S2-implícito, 0 S3. O motivo
real da reabertura foi o primeiro achado (S1); o segundo veio de carona. Pela
regra do scoreboard a rodada conta como ``reaberta por S2'' por causa do
segundo e do terceiro. Sem \emph{drip-feeding}: a revisão 2 não trouxe achado
novo. Sem \emph{churn}: a passada 2 não mexeu no código de produção.

\section{Scoreboard}

Saída de \path{sweatshop/scripts/scoreboard.ps1} sobre os quatro repositórios
desta máquina com \path{.scratch/run-log.md} (PhysiSyst, SIGAA-ME, SynchroNice,
slip), depois de registrar a atribuição acima. Para o PhysiSyst, o scoreboard
rodou sobre a cópia \path{docs/run-log/run-log.md} com as 13 linhas deste run
(154 linhas): o log local ainda não tem as 38 linhas de 30/09 da outra
máquina. Nenhum dos dois logs tem as linhas do CLEAN-23 (PR \#11).

\medskip
\noindent{\footnotesize
\begin{tabularx}{\linewidth}{@{}X r r r r r r@{}}
\toprule
Implementador & Estágios & To review & Falhou & Perguntou & Mediana até review & Custo\\
\midrule
sonnet & 121 & 75 & 31 & 15 & 7m & \$10.53 (4/121)\\
claude-opus-5-5 & 35 & 31 & 4 & 0 & 6m & \$0.72 (1/35)\\
gpt-6.1-sol high & 33 & 29 & 0 & 4 & 8m & \$23.85 (33/33)\\
sonnet-5.5 high & 8 & 8 & 0 & 0 & 6m & \$13.95 (8/8)\\
sonnet-5.5 xhigh & 2 & 2 & 0 & 0 & 19m & \$8.82 (2/2)\\
fable & 5 & 5 & 0 & 0 & 3m & ?\\
gpt-6-luna & 23 & 11 & 10 & 2 & 26m & ?\\
gpt-6-sol & 6 & 5 & 1 & 0 & 17m & ?\\
gpt-5.6-sol & 5 & 4 & 1 & 0 & 17m & ?\\
gpt-6-astra & 18 & 11 & 6 & 1 & 8m & ?\\
sonnet-5 xhigh & 11 & 6 & 2 & 3 & 8m & \$18.09 (11/11)\\
gpt-6-astra low & 2 & 2 & 0 & 0 & 30m & \$19.03 (2/2)\\
\bottomrule
\end{tabularx}}

\medskip
\noindent{\footnotesize
\begin{tabularx}{\linewidth}{@{}l X r r r r r r@{}}
\toprule
Implementador & Revisor & Reviews & Reabertos & Taxa & Por S2 & Taxa S2 & Sem rótulo\\
\midrule
sonnet & opus & 68 & 20 & 29\,\% & 1 & $\geq$1\,\% & 19\\
? & opus & 4 & 0 & 0\,\% & 0 & 0\,\% & 0\\
claude-opus-5-5 & claude-fable-5-1 & 30 & 1 & 3\,\% & 0 & $\geq$0\,\% & 1\\
gpt-6.1-sol high & gpt-6.1-sol max & 28 & 7 & 25\,\% & 6 & 21\,\% & 0\\
sonnet-5.5 high & gpt-6.1-sol max & 8 & 1 & 12\,\% & 1 & 12\,\% & 0\\
sonnet-5.5 xhigh & gpt-6.1-sol max & 2 & 1 & 50\,\% & 1 & 50\,\% & 0\\
sonnet & fable & 1 & 0 & 0\,\% & 0 & 0\,\% & 0\\
fable & fable & 4 & 1 & 25\,\% & 0 & $\geq$0\,\% & 1\\
gpt-6-luna & gpt-6-sol & 11 & 7 & 64\,\% & 0 & $\geq$0\,\% & 7\\
gpt-6-sol & gpt-5.6-sol & 5 & 4 & 80\,\% & 0 & $\geq$0\,\% & 4\\
gpt-5.6-sol & gpt-6-astra & 4 & 3 & 75\,\% & 0 & $\geq$0\,\% & 3\\
gpt-6-astra & gpt-5.6-sol & 11 & 5 & 45\,\% & 0 & $\geq$0\,\% & 5\\
sonnet-5 xhigh & opus-5.5 xhigh & 5 & 0 & 0\,\% & 0 & 0\,\% & 0\\
gpt-6-astra low & gpt-6.1-sol max & 2 & 1 & 50\,\% & 1 & 50\,\% & 0\\
\bottomrule
\end{tabularx}}

\medskip
\noindent{\footnotesize
\begin{tabularx}{\linewidth}{@{}X r r r r r r@{}}
\toprule
Implementador $\to$ Revisor & Rodadas & S2 expl.\ código & S2 expl.\ teste & S2 implícito & S1 & S3 ruído\\
\midrule
gpt-6.1-sol high $\to$ gpt-6.1-sol max & 7 & 3 & 0 & 7 & 1 & 0\\
sonnet-5.5 high $\to$ gpt-6.1-sol max & 1 & 0 & 0 & 2 & 3 & 0\\
sonnet-5.5 xhigh $\to$ gpt-6.1-sol max & 1 & 1 & 0 & 1 & 2 & 0\\
sonnet $\to$ opus & 1 & 1 & 1 & 0 & 0 & 0\\
gpt-6-astra low $\to$ gpt-6.1-sol max & 1 & 2 & 0 & 1 & 0 & 0\\
\bottomrule
\end{tabularx}}

\paragraph{O que mudou desde o relatório anterior.}
\begin{itemize}
  \item \textbf{Linha nova: \texttt{sonnet-5.5 xhigh} $\to$
        \texttt{gpt-6.1-sol max}}, 2 revisões, 1 reaberta, taxa S2 de 50\,\%.
        É deste run.
  \item \textbf{\texttt{gpt-6.1-sol high} $\to$ \texttt{gpt-6.1-sol max}} foi
        de 16 revisões e 19\,\% para 28 e 21\,\%. Este run soma 4 revisões sem
        reabertura; as outras 8 revisões e as 3 reaberturas novas vieram dos
        repositórios irmãos desde 01/10.
  \item \textbf{Linhas novas dos irmãos:} \texttt{gpt-6-astra},
        \texttt{gpt-6-sol}, \texttt{gpt-5.6-sol}, \texttt{sonnet-5 xhigh} e
        \texttt{gpt-6-astra low}. \texttt{sonnet} $\to$ \texttt{opus} ganhou 3
        revisões e uma reabertura classificada como S2. Nenhuma delas é deste
        run.
\end{itemize}

\paragraph{Como ler.}
\begin{itemize}
  \item \textbf{Mesmo revisor, três implementadores:} sob
        \texttt{gpt-6.1-sol max}, o Sol high tem taxa S2 de 21\,\% (28
        revisões), o Sonnet 5.5 high 12\,\% (8) e o Sonnet 5.5 xhigh 50\,\%
        (2). \textbf{Os dois pares do Sonnet estão abaixo de 10 revisões: são
        anedota.} O xhigh rodou um único ticket, o hard da fila.
  \item \textbf{Os 50\,\% do xhigh superestimam o implementador.} A rodada
        conta como S2 por um achado de lugar dos testes (em dúvida) e por uma
        duplicação que não bloqueou; o achado que reabriu é S1. Contando só os
        achados que bloquearam, a taxa continua em 50\,\% só por causa do segundo
        achado, que é S2 em dúvida.
  \item \textbf{5 S1 em 2 rodadas do Sonnet 5.5} apontam para o stage 1
        dos tickets que ele recebeu, PHY-45 e PHY-54, não para o
        implementador.
\end{itemize}

\section{Observações sobre os modelos}

\subsection*{GPT-6.1 Sol (high), implementação}
\begin{itemize}
  \item \textbf{4 de 4 em \texttt{to-review}, 4 de 4 mergeados de primeira.}
        6--10 min, \$0,54--0,71.
  \item \textbf{Mutate-verify por teste, com a saída vermelha copiada no
        ticket} em todos os quatro, inclusive nos testes de DOM do CLEAN-24 e
        do PHY-53, como o AGENTS.md exige.
  \item \textbf{No PHY-53, 3 dos 5 testes passaram na primeira execução}
        (ordem, ausência, limpeza): eram caracterização do comportamento que a
        primeira fatia já trazia. A prova veio das mutações, e o
        \texttt{.final} disse isso. O harness foi corrigido depois do commit
        de testes, com o vermelho repetido. A revisão aceitou.
  \item \textbf{Respeitou o escopo:} deixou a exibição do aviso para o PHY-53
        no PHY-51, como o ticket mandava, e registrou a limitação do limiar de
        1\,cm no PHY-52, prevista pelo ticket.
\end{itemize}

\subsection*{Sonnet 5.5 (xhigh), implementação hard}
\begin{itemize}
  \item \textbf{A física saiu certa na primeira passada.} Os 11 cenários de
        energia, a volta acima de $2\pi$ e as duas polias passaram; a revisão
        repetiu o vermelho (16 testes) e as cinco mutações.
  \item \textbf{Achou um defeito que o ticket não previa:} sem o clamp, uma
        corda ideal em grupo esticava de 7,66 para 8,30\,m. Escreveu o teste,
        provou com mutação e corrigiu, mas fora dos Primary files. Foi isso que
        reabriu.
  \item \textbf{Na passada 2, chamou o proxy por conta própria} para ampliar o
        contrato, em vez de remover a guarda ou escrever o
        \texttt{Proxy decided} ele mesmo. É o segundo run seguido em que o
        Sonnet 5.5 faz isso.
  \item \textbf{Caro:} \$6,06 e 110\,k de saída na passada 1, 9,8 vezes uma
        implementação do Sol high no mesmo run. Não dá para separar o efeito do
        xhigh do efeito do ticket: foi o único hard.
\end{itemize}

\subsection*{GPT-6.1 Sol (max), revisão}
\begin{itemize}
  \item \textbf{Repetiu vermelho e mutações em todas as seis revisões} e rodou
        o gate completo antes de cada merge.
  \item \textbf{Reabriu o PHY-54 pela regra, não pelo mérito:} escreveu que a
        guarda era necessária e que os critérios passavam, e mesmo assim
        devolveu, citando o \texttt{ticket-flow}. O ciclo extra custou \$3,45 e
        19 min para nenhuma linha de produção.
  \item \textbf{Achou um caso que ninguém pediu} (CLEAN-25, projeções
        invertidas), reproduzido com o código de produção, e o tirou do
        contrato em vez de reabrir por ele.
  \item \textbf{Nenhum CLEAN de documentação neste run}, ao contrário do
        anterior. A revisão do PHY-51 ajustou ela mesma a ordem de comentários.
\end{itemize}

\section{Driver e runtime}
\begin{itemize}
  \item \textbf{Limite de uso do Codex:} 00:45, ``try again at 3:49 AM''. O
        driver esperou até 03:51 e retomou sozinho, sem gastar tentativa. É o
        segundo run seguido com esse limite (o anterior às 15:44), e este
        custou 3\,h\,06 de relógio.
  \item \textbf{As recomendações do relatório anterior já estão no driver:}
        teto de revisão de 40 min, arquivo \texttt{STOP} e parada imediata num
        400. Nenhuma foi exercitada neste run: a revisão mais longa levou
        18 min.
  \item \textbf{Gate estável:} nenhum estágio precisou de
        \texttt{VITEST\_MAX\_WORKERS=1}.
  \item \textbf{Aviso inofensivo de novo:} o \texttt{.err} do Claude traz
        \texttt{claude-code:unrecognized\_model} para
        \texttt{claude-sonnet-5-5}; o estágio rodou normalmente. Os do Codex
        trazem ``failed to refresh available models: request timed out'', sem
        efeito.
  \item \textbf{Log local incompleto:} \path{.scratch/run-log.md} desta
        máquina não tem as 38 linhas de 30/09, e nenhum log tem o CLEAN-23. A
        cópia em \path{docs/run-log/} é a mais completa.
\end{itemize}

\section{Recomendações}
\begin{enumerate}
  \item \textbf{Stage 1: quando o proxy corta o escopo, listar nos Primary
        files toda função que o corte toca.} O PHY-54 nomeava
        \texttt{pieceLengths} no texto, mas não no Primary. Apoia-se nos 5 S1
        das linhas \texttt{sonnet-5.5 high} e \texttt{sonnet-5.5 xhigh}
        $\to$ \texttt{gpt-6.1-sol max} (1 rodada cada; anedota).
  \item \textbf{\texttt{ticket-flow}: o stage 3 pode consultar o proxy antes
        de reabrir só por fronteira}, quando ele mesmo julga o código fora da
        fronteira necessário. O proxy deu a resposta em minutos na passada 2; a
        reabertura custou \$3,45 e 19 min a mais.
  \item \textbf{Primary files: nomear o arquivo de teste, não o
        \texttt{describe}.} Um bloco de testes no lugar errado foi o segundo
        motivo da reabertura do PHY-54, sem efeito em comportamento.
  \item \textbf{Codex: contar com cerca de 7--8 estágios por janela} do plano
        atual em runs longos. Os dois últimos runs bateram no limite depois de
        poucas horas; um lineup com o implementador no Claude dividiria a
        carga.
  \item \textbf{Sonnet 5.5 xhigh no hard:} manter só se o próximo hard
        confirmar. Hoje há 2 revisões, e o custo (\$8,82 por ticket) é 7 vezes
        o de um ticket do Sol high (\$1,20). A recomendação não se apoia no
        scoreboard; o par tem 2 revisões.
  \item \textbf{Abrir um CLEAN para o \texttt{wrapAngle} duplicado}
        (\texttt{ropePath.ts} e \texttt{simulator.ts}). O stage 2 e a revisão
        o anotaram, mas nenhum ticket existe.
  \item \textbf{A fila acabou:} restam CLEAN-25, PHY-55 e PHY-56, todos
        \texttt{blocked} à espera do stage 1.
\end{enumerate}

\section{Débito humano}

Nenhuma linha \texttt{Débito humano:} nos cinco tickets. O PR
\href{https://github.com/Laginho/PhysiSyst/pull/12}{\#12} tem as cinco linhas
como \texttt{Review: agent} --- Approve, sem ``Needs your call''. Para o Bruno:
mergear o PR \#12, rodar o stage 1 de CLEAN-25, PHY-55 e PHY-56, e decidir
sobre o CLEAN do \texttt{wrapAngle} (recomendação 6).

\end{document}
